%%
%%  Department of Electrical, Electronic and Computer Engineering
%%  MEng Dissertation - Chapter 1: Introduction
%%

\chapter{Introduction}
\label{chap_1}

This chapter motivates multi-person pose grouping as a separable problem within bottom-up human pose estimation (HPE), identifies the gap in modular detector-agnostic grouping methods, and sets out the research questions, approach, and contributions that the rest of the dissertation defends.

\section{Context}
\label{sec_intro_context}

Introduces the detection/grouping split in bottom-up HPE: modern detectors (HigherHRNet, HRNet) produce accurate per-joint heatmaps; assigning those joints to person identities is the remaining bottleneck. Summarises why joint-to-person assignment is intrinsically hard (unknown $K$, type constraints, occlusion, scale) and why it deserves to be studied as a module in its own right.

\section{Problem Statement}
\label{sec_intro_problem}

Formulates the task: given $N$ detected keypoints with joint-type labels $t_i \in \{1,\dots,17\}$ from an image containing an unknown number of people $K$, assign each keypoint to one of $K$ person groups under the constraint that every person has at most one keypoint of each type.

\section{Research Gap}
\label{sec_intro_gap}

Existing grouping methods (associative embeddings in HigherHRNet, part-affinity fields in OpenPose, HGG, CenterGroup) are tightly coupled to their host detectors and trained jointly with them. A \emph{modular} grouping module that operates on any detector's keypoint output is largely unstudied. This is the gap this work addresses.

\section{Research Objective and Research Questions}
\label{sec_intro_objective}

The objective is to design and evaluate a detector-independent, GNN-based pose grouping module. The investigation addresses the following research questions. Bracketed short answers anticipate the findings from Chapter~\ref{chap_4}; full support is provided there.

\begin{itemize}
\item \textbf{RQ1 --- The lever.} Is the bottleneck for clustering-based pose grouping the embedding, or the grouping algorithm? \textit{[Answer: the embedding. Four learned grouping heads (DEC, slot attention, graph partitioning, SA-DMoN) fail to beat plain kNN on the same embeddings; strengthening the embedding (PG-GAT) then beats every prior configuration.]}

\item \textbf{RQ2 --- Inductive biases.} Which skeleton-aware modifications to a GATv2 embedding improve pose grouping, and by how much? \textit{[Answer: three modifications --- type-pair attention modulation, skeleton-aware positional encoding, and a type-repulsion loss term --- each contribute independently; the full combination is the best configuration.]}

\item \textbf{RQ3 --- Detector independence.} Can a grouping module trained separately from the detector match an integrated end-to-end grouping mechanism (HigherHRNet's associative embedding head) on a real-image benchmark? \textit{[Answer: not fully. PG-GAT + kNN end-to-end is approximately 0.04 PGA below HigherHRNet AE on COCO, with detector-independence as the trade-off.]}

\item \textbf{RQ4 --- Embedding geometry.} Does the choice of embedding space (Euclidean vs hyperbolic) affect pose grouping quality? \textit{[Answer: yes before COCO fine-tuning, where hyperbolic at $d=32$ matches Euclidean at $d=128$ (an efficiency edge). After fine-tuning the advantage disappears --- reported as an efficiency trade, not a quality win.]}

\item \textbf{RQ5 --- Graph primitive.} Is a higher-order graph primitive (skeleton triplets) a better node for pose grouping than single joints? \textit{[Answer: no. TriGAT underperforms joint-level PG-GAT on every comparison, with the voting step (triplet clusters to joint labels) as the dominant failure mode.]}

\item \textbf{RQ6 --- Fine-tuning.} How much does COCO fine-tuning lift a synthetically-trained embedding, and does it damage performance on the synthetic domain? \textit{[Answer: +7.5 PGA on COCO (kNN), the single largest gain of any experiment in this work; synthetic performance is not damaged (0.910 $\rightarrow$ 0.928).]}
\end{itemize}

\section{Approach}
\label{sec_intro_approach}

Train a graph attention network with contrastive loss on detected keypoints; cluster at test time with kNN. The iteration order mirrors the investigation: standard GAT baseline $\rightarrow$ a series of learned grouping heads (DEC, slot attention, graph partitioning), all of which failed to beat kNN $\rightarrow$ SA-DMoN, a skeleton-aware modularity objective, which also failed $\rightarrow$ PG-GAT, which keeps the clustering head simple but strengthens the embedding with three skeleton-aware modifications. Evaluate on synthetic data (controlled $K$, perfect GT) and COCO val2017 (zero-shot and fine-tuned). Close the loop by integrating with HigherHRNet for end-to-end validation.

\section{Research Contributions}
\label{sec_intro_contributions}

The contributions of this work are:

\begin{enumerate}
\item \textbf{PG-GAT}, a skeleton-aware graph attention network for pose-keypoint embedding. Three modifications over standard GATv2: type-pair attention modulation, skeleton-aware positional encoding, and a type-repulsion term in the contrastive loss.
\item \textbf{Empirical evidence that the embedding, not the grouping head, is the decisive lever} for clustering-based pose grouping. Documented across four learned-head approaches (DEC, slot attention, graph partitioning, and ten SA-DMoN variants), none of which outperformed plain kNN on the same embeddings.
\item \textbf{A COCO fine-tuning protocol} that lifts PG-GAT to 0.957 PGA on COCO val2017 with kNN clustering, and end-to-end to 0.942 when paired with HigherHRNet detection.
\item \textbf{An alternative-embedding-space study} (hyperbolic PG-GAT in the Lorentz model): hyperbolic at $d=32$ matches Euclidean at $d=128$ pre-fine-tune, but the advantage disappears after COCO fine-tuning --- documented as an efficiency trade, not a quality win.
\item \textbf{An alternative-primitive study} (TriGAT, triplet nodes): a negative result reported with the failure-mode analysis (voting step propagates errors).
\end{enumerate}

\section{Research Outputs}
\label{sec_intro_outputs}

The primary research output is a journal paper on PG-GAT, targeted at \emph{Image and Vision Computing} (Q1, IF 5.34), with \emph{Pattern Recognition} (Q1, IF 9.84) as a stretch target.

\section{Overview of Study}
\label{sec_intro_overview}

Chapter~\ref{chap_2} reviews the literature on bottom-up human pose estimation, graph neural networks, and clustering methods relevant to pose grouping. Chapter~\ref{chap_3} presents the mathematical preliminaries (GNN, GCN, GAT/GATv2, DEC, slot attention, graph partitioning, DMoN) and describes the methodology behind each method investigated, culminating in PG-GAT. Chapter~\ref{chap_4} reports experimental results across the methods. Chapter~\ref{chap_5} discusses what the results tell us about the problem. Chapter~\ref{chap_6} concludes and indicates future work.

\section{Declaration Regarding Use of Generative AI}
\label{sec_intro_ai}

Please declare if and how generative AI was used in this work. This must be done following the steps (adapted for dissertations/theses) outlined in 8.2.1.B.10 in the IEEE Publication Services and Products Board Operations Manual\footnote{https://pspb.ieee.org/images/files/PSPB/opsmanual.pdf}:

\begin{quote}
The use of content generated by artificial intelligence (AI) in \ldots [a dissertation/thesis] (including but not limited to text, figures, images, and code) shall be disclosed \ldots The AI system used shall be identified, and specific sections of the \ldots [dissertation/thesis] that use AI-generated content shall be identified and accompanied by a brief explanation regarding the level at which the AI system was used to generate the content.
\end{quote}

While the use of AI systems for editing and grammar enhancement is common practice, its use in this role must also be disclosed.

%% End of File.
